\section{Marco Teórico}

\subsection{Adopción de Servicios Tecnológicos en el Agro Colombiano}

El artículo de García et al. (2020) examina los determinantes de la adopción de tecnologías de la información y comunicación (TIC) en el sector agrícola colombiano, enfatizando que la mera disponibilidad de tecnología no garantiza su uso. Los autores identifican barreras económicas, como altos costos de implementación y acceso limitado a financiamiento, así como factores comportamentales, incluyendo aversión al riesgo, falta de alfabetización digital y resistencia cultural al cambio. A través de un análisis empírico basado en encuestas a productores en diversas regiones, se demuestra que las políticas de adopción tecnológica deben ir más allá de la provisión de infraestructura, incorporando estrategias de capacitación y apoyo financiero para superar estas barreras.

En el contexto del proyecto Portal Agro Comercial del Huila, este conocimiento se aplica directamente a la estrategia de implementación. El portal debe diseñarse considerando no solo la funcionalidad técnica, sino también la usabilidad para agricultores con bajo nivel de digitalización. Por ejemplo, integrar tutoriales interactivos y soporte técnico gratuito para reducir la aversión al riesgo y mejorar la alfabetización digital.

Mejoras e innovaciones implementables incluyen la integración de un sistema de incentivos económicos, como descuentos en comisiones por transacciones iniciales, o alianzas con entidades financieras para ofrecer microcréditos vinculados al uso del portal. Además, se podría adoptar un enfoque de "tecnología inclusiva" mediante interfaces multilingües y adaptadas a dispositivos móviles básicos, inspirado en modelos de adopción gradual.

Artículos similares: Rodríguez et al. (2019) sobre determinantes socioeconómicos de la adopción TIC en agricultura latinoamericana; López (2021) sobre barreras culturales en la digitalización rural.

Gráfica recomendada: Un gráfico de barras que represente la frecuencia de barreras económicas vs. comportamentales en la adopción TIC, basado en datos de encuestas a productores del Huila.

\subsection{Innovación y Marketing Agrícola en Colombia}

El trabajo de Martínez y Sánchez (2018) explora cómo la innovación en marketing agrícola puede transformar la comercialización en Colombia, destacando que la producción eficiente no es suficiente sin estrategias creativas de venta. Los autores analizan casos de éxito en el sector hortícola, donde el branding, el storytelling y el uso de redes sociales han permitido a pequeños productores diferenciarse y acceder a mercados premium. Se enfatiza la importancia de contar historias auténticas sobre el origen de los productos, promoviendo valores como sostenibilidad y calidad local, para construir lealtad del consumidor.

Aplicación al proyecto: El Portal Agro Comercial del Huila puede incorporar módulos de marketing digital, permitiendo a los productores crear perfiles con narrativas personalizadas, fotos de sus fincas y testimonios, facilitando la conexión emocional con compradores. Esto reduce la intermediación al empoderar al agricultor para vender directamente su "historia" y valor agregado.

Innovaciones: Implementar herramientas de inteligencia artificial para sugerir estrategias de marketing basadas en datos de ventas, o integrar marketplaces con funciones de e-commerce avanzadas como reseñas y calificaciones. Además, colaborar con influencers locales para amplificar el alcance.

Artículos similares: Gómez (2020) sobre marketing digital en agricultura sostenible; Pérez et al. (2022) sobre storytelling en cadenas de valor agropecuarias.

Gráfica recomendada: Un diagrama de flujo ilustrando el proceso de innovación en marketing agrícola, desde la producción hasta la venta directa, con ejemplos del Huila.

\subsection{Superando las Brechas de Transformación Digital en el Agro}

En su artículo, Torres (2019) aborda las brechas digitales en el sector agropecuario, argumentando que la transformación digital requiere no solo acceso a tecnología, sino también capacitación continua y democratización del conocimiento. Se analiza cómo las disparidades en infraestructura, educación y políticas públicas perpetúan la exclusión digital, especialmente en regiones rurales como el Huila, donde el 70\% de los productores carecen de habilidades básicas en TIC.

Aplicación: El portal debe incluir un centro de capacitación virtual con cursos gratuitos sobre uso de plataformas digitales, análisis de datos y ciberseguridad, para cerrar estas brechas y fomentar la adopción masiva.

Mejoras: Desarrollar alianzas con universidades y ONGs para programas de alfabetización digital in situ, y usar realidad aumentada para simulaciones agrícolas.

Artículos similares: Ramírez (2021) sobre inclusión digital en agricultura; Silva et al. (2023) sobre políticas de transformación digital rural.

Gráfica recomendada: Un mapa de calor mostrando brechas digitales por municipios del Huila, con datos de acceso a internet y capacitación.

\subsection{Plataformas Tecnológicas en Agricultura 4.0}

Ojeda-Beltrán (2022) describe cómo las plataformas tecnológicas sirven como columna vertebral para la Agricultura 4.0, facilitando la trazabilidad, la integración de datos y la conexión productor-mercado. El autor destaca la necesidad de plataformas interoperables que conecten sensores IoT, big data y mercados en línea, mejorando la eficiencia y reduciendo costos.

Aplicación: El Portal Agro Comercial del Huila se basa en esta premisa, actuando como plataforma central que integra datos de producción, precios y logística.

Innovaciones: Incorporar APIs para integración con sistemas de riego inteligente y drones, permitiendo trazabilidad completa desde la siembra hasta la venta.

Artículos similares: Díaz (2023) sobre plataformas IoT en agro; Mendoza (2024) sobre integración de big data en cadenas de suministro agrícolas.

Gráfica recomendada: Un diagrama de arquitectura de la plataforma, mostrando módulos de trazabilidad y conexión mercado.

\subsection{Internet de las Cosas (IoT) aplicado al Agro}

El artículo de Elijah et al. (2018) revisa aplicaciones de IoT en agricultura, enfatizando que el desafío no es la tecnología en sí, sino la conectividad rural y la interpretación de datos. Se discuten casos de monitoreo de suelos, clima y cultivos, pero se advierte sobre brechas en infraestructura.

Aplicación: Integrar sensores IoT en el portal para que productores monitoreen fincas en tiempo real, optimizando decisiones y reduciendo intermediarios al proporcionar datos verificables.

Mejoras: Usar redes LoRa para conectividad en áreas remotas del Huila, y algoritmos de machine learning para predicciones.

Artículos similares: Farooq et al. (2020) sobre IoT en agricultura de precisión; Kumar (2022) sobre desafíos de conectividad en IoT rural.

Gráfica recomendada: Un gráfico de líneas mostrando mejoras en rendimiento agrícola con IoT vs. métodos tradicionales.

\subsection{Factores de Adopción de Agricultura de Precisión}

El trabajo de Pierce y Nowak (1999) identifica factores como costos, tamaño de finca y capacitación como determinantes clave para la adopción de agricultura de precisión. Aunque clásico, su relevancia persiste en contextos colombianos.

Aplicación: El portal puede ofrecer herramientas de agricultura de precisión accesibles, como mapas geoespaciales, para fincas pequeñas.

Innovaciones: Modelos de suscripción económica para tecnologías GPS y sensores.

Artículos similares: Zhang et al. (2021) sobre adopción en países en desarrollo; Castillo (2023) sobre agricultura de precisión en Colombia.

Gráfica recomendada: Un scatter plot de costo vs. beneficio de adopción por tamaño de finca.

\subsection{E-commerce en el Sector Hortofrutícola Colombiano}

Gómez et al. (2021) analizan cómo el e-commerce reduce intermediación en horticultura, pero requiere logística y confianza. Se proponen plataformas seguras para transacciones.

Aplicación: El portal como marketplace electrónico para productos del Huila.

Mejoras: Integrar pagos digitales seguros y logística integrada.

Artículos similares: Vega (2022) sobre confianza en e-commerce agrícola; Herrera (2024) sobre logística en comercio electrónico rural.

Gráfica recomendada: Un gráfico de torta mostrando reducción de intermediarios con e-commerce.

\subsection{TIC para la Comercialización Agrícola}

Rodríguez (2017) discute cómo las TIC empoderan agricultores para mercados amplios, reduciendo dependencia de intermediarios.

Aplicación: Apps móviles en el portal para acceso directo.

Innovaciones: Chatbots para consultas de mercado.

Artículos similares: Muñoz (2019) sobre TIC en comercialización; Álvarez (2023) sobre empoderamiento digital rural.

Gráfica recomendada: Un diagrama de red mostrando conexiones productor-mercado vía TIC.

\subsection{Tendencias de Agricultura 4.0 y Tecnologías de Precisión}

Wolfert et al. (2017) revisan tendencias en Agricultura 4.0, esencial para sostenibilidad.

Aplicación: Portal como hub para tecnologías 4.0.

Mejoras: Integración con IA para recomendaciones.

Artículos similares: Klerkx et al. (2019) sobre innovación en agro; Jiménez (2024) sobre tendencias en Latam.

Gráfica recomendada: Un timeline de tendencias tecnológicas en agricultura.

\subsection{Índice de Desarrollo Digital y Competitividad}

El índice de desarrollo digital de la Unión Internacional de Telecomunicaciones (UIT, 2020) mide competitividad digital.

Aplicación: Evaluar y mejorar índice digital de productores del Huila vía portal.

Innovaciones: Dashboards de métricas digitales.

Artículos similares: Banco Mundial (2021) sobre desarrollo digital rural; Cepal (2023) sobre competitividad en agro.

Gráfica recomendada: Un radar chart comparando índices digitales por región.

\subsection{Gestión Comercial del Productor Agrícola}

Hernández (2016) propone que agricultores sean gerentes de fincas.

Aplicación: Herramientas de gestión en el portal.

Mejoras: Cursos de gestión empresarial.

Artículos similares: Soto (2020) sobre gestión en agricultura familiar; Rivera (2022) sobre emprendimiento rural.

Gráfica recomendada: Un organigrama de roles en gestión agrícola.

\subsection{Mercado de Productos Agrícolas Ecológicos en Colombia}

López y García (2018) analizan nicho ecológico con consumidores dispuestos a pagar más.

Aplicación: Sección dedicada a productos ecológicos en el portal.

Innovaciones: Certificación digital integrada.

Artículos similares: Martínez (2021) sobre consumo ecológico; Díaz (2023) sobre mercados verdes.

Gráfica recomendada: Un gráfico de barras de precios premium para productos ecológicos.

\subsection{Agricultura de Precisión con Herramientas Geoespaciales}

Auernhammer (2001) sobre dejar de cultivar a ciegas con geoespaciales.

Aplicación: Mapas interactivos en el portal.

Mejoras: Integración con satélites.

Artículos similares: McBratney et al. (2005) sobre agricultura de precisión; Vargas (2024) sobre geoespaciales en Colombia.

Gráfica recomendada: Un mapa temático de variables agrícolas.

\subsection{TIC para Fortalecer Comercialización del Agricultor (Apps y RRSS)}

Van der Burg et al. (2019) sobre smartphones como herramienta agrícola.

Aplicación: App móvil del portal con RRSS integradas.

Innovaciones: Algoritmos de matching productor-comprador.

Artículos similares: Aker (2011) sobre móviles en agricultura; Sánchez (2023) sobre RRSS en comercialización.

Gráfica recomendada: Un infográfico de uso de apps en agricultura.

\subsection{Innovación en Marketing Agrícola (Branding y Storytelling)}

Wood y O'Connell (2015) sobre vender experiencias.

Aplicación: Branding personalizado en perfiles del portal.

Mejoras: Herramientas de creación de contenido.

Artículos similares: Lindgreen et al. (2020) sobre marketing sostenible; Torres (2022) sobre storytelling digital.

Gráfica recomendada: Un storyboard de campaña de marketing agrícola.

\subsection{Agricultura 4.0: IoT, Big Data e Integración de Sistemas}

Wolfert et al. (2017) sobre sistemas integrados.

Aplicación: Portal como integrador de datos.

Innovaciones: Blockchain para trazabilidad.

Artículos similares: Kamilaris et al. (2017) sobre big data en agro; Fernández (2024) sobre integración en Colombia.

Gráfica recomendada: Un diagrama de sistemas interconectados.

\subsection{Índice Digital y Competitividad de Pequeños Productores}

UIT (2020) sobre brecha digital.

Aplicación: Medir y cerrar brechas vía portal.

Mejoras: Programas de inclusión.

Artículos similares: FAO (2021) sobre digitalización rural; Cepal (2022) sobre brechas en Latam.

Gráfica recomendada: Un gráfico de brechas digitales por grupo etario.

\subsection{Revisión Latinoamericana sobre Digitalización Agrícola}

Graef et al. (2020) sobre patrones en adopción.

Aplicación: Aprender de casos latinoamericanos para el Huila.

Innovaciones: Redes regionales.

Artículos similares: Deichmann et al. (2016) sobre digitalización en desarrollo; Rodríguez (2023) sobre Latam.

Gráfica recomendada: Un mapa regional de adopción tecnológica.

\subsection{Plataformas Tecnológicas en Agricultura 4.0 (Ojeda-Beltrán, 2022)}

Ya mencionado, enfocar en conexión con necesidades reales.

Aplicación: Portal adaptado a productores del Huila.

Mejoras: Feedback loops para mejoras.

Artículos similares: Ojeda-Beltrán (2023) extensiones; López (2024) sobre plataformas locales.

Gráfica recomendada: Un flowchart de interacción usuario-plataforma.

Análisis comparativo: Los autores como García y Ojeda coinciden en barreras económicas, mientras Torres enfatiza capacitación. Elijah y Wolfert destacan integración tecnológica, diferenciándose de enfoques comportamentales de Martínez.